%% Solutions for Chapter 11: Turing Machines

\chapter{Turing Machines --- Solutions}

\begin{enumerate}[{11.}1. ]

%%------------------------------------------------------------
\item \textbf{How is a Turing machine like an elevator?}

\textbf{Analogy}:
\begin{itemize}
\item \textbf{States} correspond to \emph{floors} (the elevator's current position/state).
\item \textbf{Input} corresponds to \emph{button presses} (the sequence of requested floors).
\item The elevator transitions from floor to floor based on button requests, just as the Turing machine transitions from state to state based on input symbols.
\end{itemize}

\textbf{Essential missing component}: An elevator has only \emph{finite, bounded memory} --- it can remember which buttons have been pressed (a finite set), but it cannot write new symbols to an unbounded tape. The critical component preventing an elevator from modeling a general computing device is the \emph{unbounded read/write tape}. Without it, the elevator can only implement a finite state machine (regular language computation), not a Turing machine. The ability to store and retrieve arbitrarily large amounts of information is essential for general computation.

%%------------------------------------------------------------
\item \textbf{Turing machine for palindromes $L = \{w \mid w = w^r\}$ over $\Sigma=\{\pmb{a},\pmb{b},\pmb{c}\}$.}

\begin{enumerate}
\item \textbf{One-tape, one-head TM}:

Strategy: Repeatedly match the leftmost and rightmost unprocessed symbols.
\begin{enumerate}[1.]
\item Mark the leftmost symbol $a$ (replace with $\overline{a}$); move right to the rightmost unmarked symbol.
\item If the rightmost symbol matches $a$, mark it and move left; continue.
\item If the rightmost symbol does not match, reject.
\item When all symbols are marked (or only the center symbol for odd-length), accept.
\end{enumerate}
States: $\{q_\text{start}, q_a, q_b, q_c, q_\text{right}, q_{\text{check},a}, q_{\text{check},b}, q_{\text{check},c}, q_\text{return}, q_\text{accept}, q_\text{reject}\}$.

Transitions: $\delta(q_\text{start}, \pmb{a}) = (q_a, \overline{\pmb{a}}, R)$ [mark $\pmb{a}$, go right]; etc.

\item \textbf{Linear Bounded Automaton (LBA)}: Since the tape head never needs to go beyond the input boundaries (matching leftmost/rightmost within the input), the TM above is already an LBA. The computation uses only the $n$ tape cells containing the input.

\item \textbf{Nondeterminism/multiple tapes}: A two-tape TM copies the input to tape 2, reverses tape 2, and compares tape 1 with tape 2 symbol by symbol --- $O(n)$ time. A nondeterministic TM can guess the center and verify matching from both ends simultaneously.
\end{enumerate}

%%------------------------------------------------------------
\item \textbf{TM recognizing $\{\pmb{a}^n \mid n \text{ is a perfect square}\}$ over $\Sigma=\{\pmb{a}\}$.}

\textbf{Strategy}: The $k$-th perfect square is $k^2$. Note $(k+1)^2 = k^2 + (2k+1)$: successive perfect squares differ by successive odd numbers $1, 3, 5, 7, \ldots$

Algorithm:
\begin{enumerate}[1.]
\item If input is $\lambda$: accept (since $0 = 0^2$).
\item Mark the first $\pmb{a}$ and subtract $1$; this corresponds to $1^2=1$.
\item Then subtract $3$; if exhausted, accept (that was $2^2=4$). If not enough, reject.
\item Then subtract $5$; if exhausted accept ($3^2=9$). Etc.
\item Subtract $2k+1$ at each step for $k=0,1,2,\ldots$; accept if the tape is exactly emptied, reject if running short.
\end{enumerate}

The TM maintains the current step $k$ (in unary, using a separate tape region or using a scratch area) and repeats the subtraction of $2k+1$. It accepts when exactly $k^2$ symbols have been processed.

%%------------------------------------------------------------
\item \textbf{TM for languages over $\Sigma=\{\pmb{a},\pmb{b},\pmb{c}\}$.}

\begin{enumerate}
\item $\{\pmb{a}^k\pmb{b}^n\pmb{c}^m \mid k\neq n \wedge n\neq m\}$:

Strategy: Count $|x|_a$, $|x|_b$, $|x|_c$ separately and verify $|x|_a \neq |x|_b$ and $|x|_b \neq |x|_c$.

A practical approach:
\begin{enumerate}[1.]
\item Verify the input has the form $\pmb{a}^k\pmb{b}^n\pmb{c}^m$ (check symbol ordering); reject if not.
\item Pair off $\pmb{a}$s and $\pmb{b}$s: erase one $\pmb{a}$ and one $\pmb{b}$ in each pass. If they run out simultaneously, $k=n$; reject.
\item Pair off $\pmb{b}$s and $\pmb{c}$s similarly. If they run out simultaneously, $n=m$; reject.
\item If both checks show inequality, accept.
\end{enumerate}

\item $\{x \mid |x|_a \neq |x|_b \wedge |x|_b \neq |x|_c\}$ (arbitrary mixing):

Strategy: use two separate matching rounds, with fresh marks in each round, so
that the $\pmb{b}$'s are still available for the second comparison.
\begin{enumerate}[1.]
\item In round 1, repeatedly mark one unmarked $\pmb{a}$ as $A$ and one unmarked
$\pmb{b}$ as $B$. If one kind runs out before the other, then
$|x|_{\pmb{a}}\neq |x|_{\pmb{b}}$ is confirmed. If they run out together, reject.
\item Restore all $B$'s to ordinary $\pmb{b}$'s.
\item In round 2, repeatedly mark one unmarked $\pmb{b}$ as $B$ and one
unmarked $\pmb{c}$ as $C$. If one kind runs out before the other, then
$|x|_{\pmb{b}}\neq |x|_{\pmb{c}}$ is confirmed. If they run out together, reject.
\item Accept iff both checks confirm inequality.
\end{enumerate}
\end{enumerate}

%%------------------------------------------------------------
\item \textbf{Machine $M$ with $L(M)=L_1(M)=L_2(M)=L_3(M)=\{x\in\{\pmb{a},\pmb{b},\pmb{c}\}^*\mid |x|_a=|x|_b=|x|_c\}$.}

\begin{enumerate}
\item \textbf{Construction}: First let $M$ deterministically test whether the input has equally many $\pmb{a}$'s, $\pmb{b}$'s, and $\pmb{c}$'s by the usual crossing-off routine. Then:
\begin{enumerate}[1.]
\item If the test succeeds, erase the whole tape, write a single $\pmb{Y}$, position the head on that $\pmb{Y}$, and halt in a final state.
\item If the test fails, do \emph{not} halt: enter a simple two-state loop that moves right and left forever on blanks.
\end{enumerate}
For a yes-instance, the machine halts in a final state, with $\pmb{Y}$ somewhere on the tape and in fact as the rightmost nonblank symbol. Thus the input belongs to all four sets $L(M),L_1(M),L_2(M),L_3(M)$. For a no-instance, the machine never halts, so it belongs to none of them. Therefore all four acceptance notions define exactly the required language.

\item \textbf{Is this always possible?}: Yes, for \emph{every} Turing-acceptable
language $L$. Let $A$ be any Turing machine with $L(A)=L$. Build a new machine
$M$ that simulates $A$, but modifies only the accepting behavior:
\begin{enumerate}[1.]
\item If $A$ reaches an accepting halt, then $M$ erases the tape, writes a single
$\pmb{Y}$, positions the head on that $\pmb{Y}$, enters a final state, and halts.
\item If $A$ would reject, or if $A$ never accepts, then $M$ simply never halts
(for example, it can enter the same two-state blank-tape loop used in part~(a)).
\end{enumerate}
Then $x\in L$ exactly when $M$ halts, exactly when it halts with $\pmb{Y}$ on the
tape, exactly when it halts with $\pmb{Y}$ at the right end of the tape, and
exactly when it halts in a final state. Hence
\[
L(M)=L_1(M)=L_2(M)=L_3(M)=L.
\]
\end{enumerate}

%%------------------------------------------------------------
\item \textbf{$L = \{ww \mid w\in\{\pmb{a},\pmb{b},\pmb{c}\}^*\}$.}

\begin{enumerate}
\item \textbf{One-tape TM}:
Strategy:
\begin{enumerate}[1.]
\item Sweep right across the input and write its length $n$ in binary in a scratch
region of the same tape to the right of the word. If $n$ is odd, reject.
\item Replace that counter by $m=n/2$.
\item For each $i=1,2,\ldots,m$, use a second binary counter on the scratch region
to locate the $i$th symbol of the word and remember whether it is $\pmb{a}$,
$\pmb{b}$, or $\pmb{c}$.
\item Using the counters again, locate the $(m+i)$th symbol and compare it with
the remembered symbol. If they differ, reject.
\item If all $m$ comparisons succeed, accept.
\end{enumerate}

\item \textbf{LBA}: Since the computation uses only the $2n$ cells of the input, the machine is already an LBA.

\item \textbf{Multiple tapes}: A two-tape TM copies the first half of the input to tape 2, then compares tape 2 with the second half of tape 1 symbol by symbol. A nondeterministic TM can guess the midpoint and verify.
\end{enumerate}

%%------------------------------------------------------------
\item \textbf{Lexicographic ordering Turing machines.}

Given $\Sigma=\{\pmb{a}_1,\ldots,\pmb{a}_n\}$ with lexicographic order based on base-$n$ numbering.

\begin{enumerate}
\item \textbf{Successor TM}: Machine $M_{succ}$ maps each word to the next word in
the bijective base-$n$ ordering described in the exercise.

Algorithm: If the tape is blank (that is, the word is $\lambda$), write
$\pmb{a}_1$. Otherwise process the word from right to left. Find the rightmost
symbol $\pmb{a}_i$ with $i<n$; replace it with $\pmb{a}_{i+1}$ and replace every
symbol to its right with $\pmb{a}_1$. If every symbol is $\pmb{a}_n$, so the word
is $\pmb{a}_n^k$, replace it by $\pmb{a}_1^{k+1}$. This is exactly the next word in
the stated order.

\item \textbf{Generator TM}: The initial blank tape already represents the first
word $\lambda$. Apply $M_{succ}$ to obtain $\pmb{a}_1$, then repeatedly erase the
current word and replace it by its successor. Thus the visible word sequence is
\[
\lambda,\ \pmb{a}_1,\ \pmb{a}_2,\ \ldots,\ \pmb{a}_n,\ \pmb{a}_1\pmb{a}_1,\ldots
\]
in lexicographic order.

\item \textbf{Enumerator TM}: Again the blank tape represents the first word
$\lambda$. Keep the current word in the rightmost block of nonblank cells. To
append the next word, copy that rightmost block to a fresh work area, run
$M_{succ}$ on the copy, and leave the result to the right of one blank separator.
Repeating this appends each successive word while leaving the earlier ones
untouched.

\item \textbf{Deterministic simulation of NDTM}: Given a nondeterministic TM $N$, simulate it deterministically using the lexicographic ordering of computation paths. At step $k$, enumerate all possible $k$-step computation paths (sequences of transitions) using the lexicographic order, and simulate each one. If any path reaches an accepting state, accept. This gives a deterministic simulation (though potentially much slower). $M_{succ}$ generates the next computation path to try.
\end{enumerate}

%%------------------------------------------------------------
\item \textbf{Semi-infinite tape Turing machine.}

A \emph{semi-infinite tape} has a left boundary (leftmost cell = cell 1) and extends indefinitely to the right.

\begin{enumerate}
\item \textbf{Two-track simulation}: Store the nonnegative cells
$0,1,2,\ldots$ of the old tape on the upper track and the negative cells
$-1,-2,-3,\ldots$ on the lower track in reverse order. Thus column $0$ contains
the old cell $0$ on the upper track and a special boundary marker on the lower
track, while column $k\geq 1$ contains old cell $k$ on the upper track and old
cell $-k$ on the lower track.

The finite control remembers whether the simulated head is on the upper or lower
track. Moving right on the old nonnegative side means moving one column right on
the upper track; moving left from old cell $0$ switches to the lower track in
column $1$; moving right from old cell $-1$ switches back to the upper track in
column $0$; and all other moves on the negative side are simulated by moving one
column left or right on the lower track as appropriate.

\item \textbf{Equivalence proof}: Induct on the number of simulated steps. The
induction invariant is that the two tracks encode exactly the same tape contents
as the original bi-infinite tape, and that the simulated head position and state
are correctly represented by the current column together with the remembered
track. Each move of the old machine changes only one scanned cell and shifts the
head by one position, and the semi-infinite simulation reproduces exactly that
effect. Therefore the two machines accept and reject the same words. \qed
\end{enumerate}

%%------------------------------------------------------------
\item \textbf{TM recognizing $\{\pmb{a}^n \mid n \text{ is a power of 2}\}$.}

Strategy: Repeatedly check if $n$ is divisible by 2 and halve it until reaching 1.

Algorithm:
\begin{enumerate}[1.]
\item If input is $\pmb{a}$ ($n=1=2^0$), accept.
\item If input has odd length $> 1$, reject.
\item Sweep left to right, marking every other $\pmb{a}$ with $\overline{\pmb{a}}$. If $n$ is odd (and $>1$) an unmarked $\pmb{a}$ remains at the end — reject. Otherwise, replace all unmarked $\pmb{a}$s with blanks; the $\overline{\pmb{a}}$s form the halved string. Repeat.
\item Repeat until one symbol remains; accept.
\end{enumerate}

%%------------------------------------------------------------
\item \textbf{Three-head TM for $\{x \mid |x|_a = |x|_b = |x|_c\}$.}

\textbf{Design}: All three heads start at the leftmost character.

Strategy:
- Head 1 scans for $\pmb{a}$-symbols (moving right, skipping $\pmb{b}$ and $\pmb{c}$).
- Head 2 scans for $\pmb{b}$-symbols.
- Head 3 scans for $\pmb{c}$-symbols.

In each synchronization step: advance all three heads to their next unmatched symbol of the respective types. If at any point one head reaches the end while others haven't (or vice versa), reject. If all three heads reach the end simultaneously, accept.

\textbf{No head ever needs to move left}: Each head only moves right (it scans past symbols of other types, looking for its own). Since the heads never backtrack, none needs to move left. \qed

%%------------------------------------------------------------
\item \textbf{Every CFL is Turing-acceptable.}

\textbf{Proof} (Lemma 11.1): Let $L$ be a CFL with PDA $P = \langle\Sigma,\Gamma,S,s_0,\delta,Z_0,F\rangle$.

Construct a TM $M$ that simulates $P$:
- Tape 1: input string $x$.
- Tape 2: stack of $P$ (top of stack at left).
- $M$ simulates each PDA move: reads next input symbol (or $\lambda$), looks up $\delta$, pops the stack top, pushes the new stack string (nondeterministically if $|\delta|>1$).
- If the PDA reaches a final state with empty input, $M$ accepts.

Since the PDA is nondeterministic, $M$ is a nondeterministic TM. Every NTM can be simulated by a deterministic TM (by enumerating computation paths as in Exercise 11.7(d)).

Thus $L = L(P) = L(M) \in \TSIG$. \qed

%%------------------------------------------------------------
\item \textbf{Every type 1 language is a LBL.}

\textbf{Proof} (Theorem 11.3): Let $G=\langle\Sigma,\Omega,Z,P\rangle$ be a
context-sensitive grammar. We follow the construction discussed just before
Theorem~11.3, namely the two-tape machine used in the proof of
Theorem~11.2.

On input $x$ of length $n$, tape 1 holds $x$ and is never altered. Tape 2 starts
with just the start symbol $Z$. The finite control nondeterministically chooses
a production from $P$ and a position in the current sentential form on tape 2,
and applies that production exactly as in Theorem~11.2.

The only extra observation needed for a type~1 grammar is that, once the machine
has begun deriving a nonempty word, the length of the sentential form can never
decrease. Therefore any branch on which tape 2 ever becomes longer than $n$ can
be abandoned immediately, because it can never later shrink back to the target
word $x$. Thus every surviving computation keeps tape 2 within $n$ cells.

If $x=\lambda$, the only successful derivation is the special initial production
$Z\to\lambda$. If $x\neq\lambda$, the machine accepts exactly when some sequence
of guessed productions makes tape 2 agree with tape 1. So the same construction
is now bounded by a linear amount of tape, and hence is a nondeterministic LBA.
Therefore every type~1 language is a LBL. \qed

%%------------------------------------------------------------
\item \textbf{Convert LBA $M$ to three-track strict LBA.}

\begin{enumerate}
\item \textbf{New alphabets}:
- Track 1 carries the old tape alphabet $\Gamma$, not merely the input alphabet
$\Sigma$.
- Track 2 marks the left end of the original word.
- Track 3 marks the right end of the original word.

It is convenient to compress these three tracks into a single alphabet
\[
\Gamma'=\Gamma\cup\{\overline{d},\underline{d}\mid d\in\Gamma\},
\]
where $\overline{d}$ means that $d$ is currently in the leftmost cell of the
word and $\underline{d}$ means that $d$ is currently in the rightmost cell. For
words of length at least $2$, these two markers occur in different cells.

\item \textbf{New transitions}: For any old transition $\delta(s,a) = (s',a',d)$ in $M$:
  - If the old move stays inside the word, copy it directly, preserving any
  overbar or underline markers on the end cells.
  - If the old machine would move left off the word to scan the left delimiter,
  replace that excursion by a short detour through a new state that records
  ``the left delimiter is being scanned.'' Because the delimiter can only return
  to the word in one way, this detour is unique.
  - Do the symmetric construction on the right end of the word for moves that
  would have scanned the right delimiter.

In other words, the strict machine does \emph{not} reject when the old machine
would have stepped onto a delimiter; instead, it simulates the unique
delimiter-scan-and-return behavior without ever leaving the marked end cells.

\item \textbf{Proof of equivalence} (for $|x|\geq 2$): By induction on the number of steps:
- Every ordinary move of $M$ is copied verbatim by $M'$.
- Every two-step excursion of $M$ from an end cell onto a delimiter and back is
replaced by the corresponding unique boundary detour in $M'$.
- Hence the two machines visit corresponding configurations inside the word and
accept exactly the same words.

So for every input of length at least $2$, the strict three-track machine
accepts exactly when the original LBA does. \qed
\end{enumerate}

%%------------------------------------------------------------
\item \textbf{Handle strings of length 1; strict LBA for type 1 languages.}

\begin{enumerate}
\item \textbf{Strings of length 1}: For each tape symbol $d$, add a special symbol
$\overline{\underline{d}}$ meaning ``this cell is both the leftmost and the
rightmost cell.'' The modified strict LBA treats these symbols as the obvious
single-cell analogues of $\overline{d}$ and $\underline{d}$.

\item \textbf{Strict LBA accepting $\lambda$}: If the initial scan is of a blank
cell, interpret that as the empty input and have the machine decide immediately
whether $\lambda$ should be accepted. After that special start-up case, the
strict machine proceeds exactly as in Exercise~11.13 and never actively moves to
a cell outside the input. Thus every type~1 language can be recognized by a
strict LBA in this slightly relaxed sense. \qed
\end{enumerate}

%%------------------------------------------------------------
\item \textbf{Prove (Theorem 11.4) by induction: TM computation iff LBA grammar derivation.}

\textbf{Induction statement}: For every $k\geq 0$,
\[
\alpha s\beta \Dvdash^k \gamma t\omega
\]
if and only if there exist $i,j,m,n\in\NN$ such that
\[
[\#^i\alpha s\beta\#^j] \DRightarrow^k [\#^m\gamma t\omega\#^n],
\]
where the derivation uses only class~2 productions.

\textbf{Base case} ($k=0$): $\alpha s\beta \Dvdash^0 \alpha s\beta$, and likewise
\[
[\#^i\alpha s\beta\#^j] \DRightarrow^0 [\#^i\alpha s\beta\#^j].
\]

\textbf{Induction step}: Assume the statement holds for $k$. Suppose
\[
\alpha s\beta \Dvdash^{k+1} \mu u\nu.
\]
Then there is an intermediate configuration $\gamma t\omega$ with
\[
\alpha s\beta \Dvdash^k \gamma t\omega \Dvdash \mu u\nu.
\]
By the induction hypothesis,
\[
[\#^i\alpha s\beta\#^j] \DRightarrow^k [\#^m\gamma t\omega\#^n]
\]
for suitable $i,j,m,n$. The final Turing-machine move is produced by one
transition of $M$, and Definition~11.8 includes exactly one corresponding class~2
production. Applying that production yields
\[
[\#^m\gamma t\omega\#^n] \DRightarrow [\#^p\mu u\nu\#^q]
\]
for some $p,q\in\NN$, depending only on whether the head has just moved onto a
new blank cell. Hence
\[
[\#^i\alpha s\beta\#^j] \DRightarrow^{k+1} [\#^p\mu u\nu\#^q].
\]

The converse is identical: every class~2 production encodes one legal move of
$M$, so a derivation of length $k+1$ factors into a derivation of length $k$
followed by one TM move. Therefore the claim holds for all $k$. \qed

%%------------------------------------------------------------
\item \textbf{General induction statement for Theorem 11.5.}

\textbf{Theorem 11.5} establishes that every LBL is a type~1 language.

\textbf{General induction statement}: Fix a strict LBA $B$ and a nonempty input
word of length at least $2$, written $\pmb{a}x\pmb{b}$. For each configuration
$C$ of $B$, let $\phi(C)$ be the corresponding sentential form in which the
scanned tape symbol carries the current state as a subscript, while the first
and last tape symbols retain their overbar and underline markers. Then for every
$k\geq 0$ and every configuration $C$,
\[
s_0\pmb{a}x\pmb{b} \Dvdash^k C
\quad\text{iff}\quad
\overline{\pmb{a}}_0x\underline{\pmb{b}} \DRightarrow^k \phi(C),
\]
where only class~2 productions are used.

\textbf{Base case} ($k=0$): The initial configuration is exactly
$s_0\pmb{a}x\pmb{b}$, and its encoding is
$\overline{\pmb{a}}_0x\underline{\pmb{b}}$.

\textbf{Induction step}: Assume the claim for $k$. If
\[
s_0\pmb{a}x\pmb{b} \Dvdash^{k+1} C',
\]
choose $C$ with
\[
s_0\pmb{a}x\pmb{b} \Dvdash^k C \Dvdash C'.
\]
By the induction hypothesis,
\[
\overline{\pmb{a}}_0x\underline{\pmb{b}} \DRightarrow^k \phi(C).
\]
Definition~11.9 gives one class~2 production for each possible move of the
strict LBA, including the boundary cases where the scanned symbol is marked with
an overbar or underline. Hence the last move $C\Dvdash C'$ is mirrored by one
derivation step $\phi(C)\DRightarrow\phi(C')$. The converse direction is the
same, since every class~2 production encodes a legal move of $B$ and nothing
else. This is the induction statement needed in the proof of
Theorem~11.5. \qed

%%------------------------------------------------------------
\item \textbf{Type 0 language Kleene closure.}

\begin{enumerate}
\item \textbf{Flaw in $G_* = \langle\Sigma,\Omega\cup\{W\},W,P\cup\{W\to\lambda,W\to WW,W\to Z\}\rangle$}:

The difficulty is that the two occurrences of $Z$ generated by $W\to WW$ are
placed \emph{adjacent} to one another. Since the productions of a type~0 grammar
may inspect several neighboring symbols at once, a production from $P$ can then
act across the boundary between the two intended copies. So the construction can
generate strings that are not concatenations of words from $L(G)$.

\item \textbf{Example}: Let
\[
G=\langle\{a,b\},\{Z,A,B\},Z,\{Z\to AB,\ A\to\lambda,\ B\to b,\ BA\to a\}\rangle.
\]
For a single copy of $Z$, the rule $BA\to a$ is unusable, so $L(G)=\{b\}$. But in
the naive construction,
\[
W\Rightarrow WW\Rightarrow ZZ\Rightarrow ABAB,
\]
and now the middle substring $BA$ straddles the boundary between the two copies.
Applying $BA\to a$ there gives
\[
ABAB\Rightarrow AaB\Rightarrow ab.
\]
Since $ab\notin\{b\}^*$, the naive construction is wrong.

\item \textbf{Correct $G_*$}: Insert a separator that does not occur in any
left-hand side of a rule from $P$. Choose a new nonterminal $\#$ with
$\#\notin\Gamma\cup\Sigma$, and define
\[
G_*=\langle\Sigma,\Gamma\cup\{W,\#\},W,P_*\rangle,
\]
where
\[
P_* = P\cup\{W\to\lambda,\ W\to Z,\ W\to Z\#W,\ \#\to\lambda\}.
\]
The symbol $\#$ separates one copy of $G$ from the next, so productions from $P$
cannot straddle two copies.

\item \textbf{Proof}: If
\[
w=w_1w_2\cdots w_k\qquad\text{with each }w_i\in L(G),
\]
then
\[
W\Rightarrow Z\#W\Rightarrow Z\#Z\#W\Rightarrow\cdots\Rightarrow
Z\#Z\#\cdots\#Z
\]
with $k$ copies of $Z$, each copy derives its own $w_i$ by the productions in
$P$, and finally all separators disappear by repeated use of $\#\to\lambda$.
Hence $w\in L(G_*)$.

Conversely, any derivation in $G_*$ first creates a string
$Z\#Z\#\cdots\#Z$, and the only productions that can affect symbols between two
occurrences of $\#$ are the productions from $P$ applied within that block, or
the final erasures of $\#$. Since no left-hand side from $P$ contains $\#$, the
copies do not interact. Therefore every terminal word derived from $W$ is a
concatenation of finitely many words from $L(G)$. So
\[
L(G_*)=L(G)^*.
\]
\qed
\end{enumerate}

%%------------------------------------------------------------
\item \textbf{\TSIG\ is closed under various operations.}

\begin{enumerate}
\item \textbf{Homomorphism}: Let $M$ accept $L$, and let $h:\Sigma\to\Gamma^*$ be
a homomorphism. On input $y\in\Gamma^*$, enumerate the words of $\Sigma^*$ in
lexicographic order (Exercise~11.7), and dovetail the simulations of $M$ on
those words. Whenever one simulation accepts some $x$, compute $h(x)$ and
compare it with $y$; if $h(x)=y$, accept. If $y\in h(L)$, some $x\in L$ with
$h(x)=y$ will eventually be found, so $h(L)\in\TSIG$.

\item \textbf{Inverse homomorphism}: On input $x$, compute $h(x)$ and simulate
$M$ on $h(x)$. Then the new machine accepts exactly the words in
$h^{-1}(L)$.

\item \textbf{Concatenation}: Let $M_1$ and $M_2$ accept $L_1$ and $L_2$. On
input $y$, list all splits $y=uv$. There are only finitely many. Dovetail the
simulations of $M_1$ on each $u$ and $M_2$ on the corresponding $v$, and accept
as soon as one pair both accepts. This semidecides $L_1L_2$.

\item \textbf{Substitution}: Let $s$ be a substitution with $s(\pmb{a})\in\TSIG$
for each $\pmb{a}\in\Sigma$, and let $M$ accept $L$. On input $y$, enumerate all
pairs consisting of a word $x=\pmb{a}_1\cdots\pmb{a}_k\in\Sigma^*$ and a
factorization $y=y_1\cdots y_k$ (allowing some $y_i=\lambda$). For each such
pair, dovetail the computation of $M$ on $x$ together with the computations of
machines for $s(\pmb{a}_i)$ on $y_i$. Accept when one pair makes all these
machines accept. That happens exactly when $y\in s(L)$. Hence \TSIG\ is closed
under substitution. \qed
\end{enumerate}

%%------------------------------------------------------------
\item \textbf{Equivalence of the acceptance modes $L_1$, $L_2$, and $L_3$.}

\begin{enumerate}
\item \textbf{$L_1(A_1)$ to $L_2(A_2)$}: Let
$A_1=\langle\Sigma,\Gamma,S,s_0,\delta\rangle$ and suppose $L=L_1(A_1)$. Form
$A_2$ by adjoining a short clean-up routine with new states $q_L,q_R,h'$. Every
old transition that formerly entered the halt state $h$ is redirected to $q_L$.
In $q_L$, the machine sweeps to the left end of the nonblank portion of the
tape; in $q_R$, it erases everything while moving right until the first blank;
then it writes a single $\pmb{Y}$ and enters $h'$. Thus $A_2$ halts with
$\pmb{Y}$ on the tape exactly on the inputs for which $A_1$ halts at all. Hence
$L_2(A_2)=L_1(A_1)=L$.

\item \textbf{$L_2(A_2)$ to $L_3(A_3)$}: Let $A_2$ be a machine with
$L=L_2(A_2)$. Again redirect each old transition to the halt state into a new
inspection state. That routine scans the tape once to determine whether any
$\pmb{Y}$ occurs. If one is found, it erases the tape, writes a single
$\pmb{Y}$, and halts. If no $\pmb{Y}$ is found, it halts with some other final
tape contents, for example a single $\pmb{N}$. Therefore $A_3$ halts with
$\pmb{Y}$ at the right end of the tape exactly when $A_2$ would have halted with
$\pmb{Y}$ somewhere on the tape. So $L_3(A_3)=L_2(A_2)=L$. \qed
\end{enumerate}

%%------------------------------------------------------------
\item \textbf{\OSIG\ (context-sensitive languages) is closed under the same operations.}

For a fixed input, an LBA has only finitely many configurations, so any LBA can
be modified to halt by rejecting as soon as a configuration repeats. Also, for
context-sensitive languages the standard closure statements for homomorphism and
substitution are understood in the \emph{nonerasing} case. That is the
interpretation used in parts (a) and (d).

\begin{enumerate}
\item \textbf{Homomorphism}: Let $h:\Sigma\to\Gamma^+$ be nonerasing and let
$L\in\OSIG$. By Corollary~11.2 there is a halting LBA $B$ for $L$. On input
$y\in\Gamma^*$, enumerate all $x\in\Sigma^*$ with $h(x)=y$. Because $h$ is
nonerasing, every such $x$ has length at most $|y|$, so only finitely many words
need be checked. Run $B$ on each of them and accept iff one is accepted. This
halts on every input, so $h(L)\in\OSIG$.

\item \textbf{Inverse homomorphism}: On input $x$, compute $h(x)$ and run the
halting machine for $L$ on that word. Therefore $h^{-1}(L)\in\OSIG$.

\item \textbf{Concatenation}: Given halting machines for $L_1$ and $L_2$, list all
splits $y=uv$ of the input and test each one. Since the list of splits is finite
and both machines halt on every input, this decides $L_1L_2$.

\item \textbf{Substitution}: Let $s$ be a nonerasing substitution with
$s(\pmb{a})\in\OSIG$ for each $\pmb{a}\in\Sigma$. On input $y$, enumerate all
factorizations $y=y_1\cdots y_k$ into nonempty pieces and all words
$x=\pmb{a}_1\cdots\pmb{a}_k\in\Sigma^k$ of the same length $k$. Because the pieces are
nonempty, only finitely many such choices exist. Test whether $x\in L$ and
$y_i\in s(\pmb{a}_i)$ for every $i$; accept iff one choice works. Hence
\OSIG\ is closed under nonerasing substitution. \qed
\end{enumerate}

\end{enumerate}
